% Created 2024-12-10 mar 16:09
% Intended LaTeX compiler: pdflatex
\documentclass[presentation]{beamer}
\usepackage[utf8]{inputenc}
\usepackage[T1]{fontenc}
\usepackage{graphicx}
\usepackage{longtable}
\usepackage{wrapfig}
\usepackage{rotating}
\usepackage[normalem]{ulem}
\usepackage{amsmath}
\usepackage{amssymb}
\usepackage{capt-of}
\usepackage{hyperref}
\usepackage{minted}
\usetheme{metropolis}
\usetheme{default}
\author{José Manuel Sánchez Álvarez}
\date{\today{}}
\title{Sistemas de Almacenamiento Distribuido}
\hypersetup{
 pdfauthor={José Manuel Sánchez Álvarez},
 pdftitle={Sistemas de Almacenamiento Distribuido},
 pdfkeywords={},
 pdfsubject={},
 pdfcreator={Emacs 27.1 (Org mode 9.6.18)}, 
 pdflang={English}}
\begin{document}

\maketitle

\begin{frame}[label={sec:orgc011e4e}]{Introducción}
El almacenamiento distribuido es clave para gestionar datos en Big Data.
\begin{itemize}
\item Divide datos entre múltiples nodos.
\item Beneficios:
\begin{itemize}
\item Escalabilidad.
\item Tolerancia a fallos.
\item Rendimiento eficiente.
\item Durabilidad.
\end{itemize}
\end{itemize}

---
\end{frame}

\begin{frame}[label={sec:org7d37574}]{¿Qué es el almacenamiento distribuido?}
\begin{block}{Definición}
\begin{itemize}
\item Divide los datos en fragmentos y los almacena en múltiples máquinas que simulan un único sistema.
\item Ventajas:
\begin{enumerate}
\item \alert{\alert{Escalabilidad horizontal:}} Agrega nodos para manejar el crecimiento de datos.
\item \alert{\alert{Tolerancia a fallos:}} Datos accesibles desde nodos redundantes.
\item \alert{\alert{Eficiencia:}} Reduce cuellos de botella.
\item \alert{\alert{Durabilidad:}} Resiliencia ante fallos físicos.
\end{enumerate}
\end{itemize}

---
\end{block}
\end{frame}

\begin{frame}[label={sec:org317cb5a}]{Sistemas Principales}
\begin{block}{HDFS (Hadoop Distributed File System)}
\begin{itemize}
\item Parte del ecosistema Apache Hadoop.
\item Diseñado para grandes volúmenes de datos.
\end{itemize}

\begin{block}{Características Clave}
\begin{itemize}
\item Alto rendimiento para operaciones secuenciales.
\item Escalabilidad horizontal.
\item Replicación de bloques (por defecto 3 veces).
\item Integración con herramientas como Apache Hive y Spark.
\end{itemize}
\end{block}

\begin{block}{Arquitectura}
\begin{enumerate}
\item \alert{\alert{NameNode:}} Administra metadatos.
\item \alert{\alert{DataNodes:}} Almacenan bloques de datos.
\end{enumerate}
\end{block}

\begin{block}{Casos de Uso}
\begin{itemize}
\item Procesamiento analítico masivo.
\item Gestión de datos históricos.
\end{itemize}

---
\end{block}
\end{block}

\begin{block}{Amazon S3}
\begin{itemize}
\item Servicio de almacenamiento de objetos en la nube de AWS.
\end{itemize}

\begin{block}{Características Clave}
\begin{itemize}
\item Almacenamiento basado en objetos con metadatos únicos.
\item Escalabilidad automática.
\item Durabilidad de \alert{\alert{11 nueves}}.
\item Integración con AWS Lambda, Athena, y más.
\end{itemize}
\end{block}

\begin{block}{Clases de Almacenamiento}
\begin{enumerate}
\item \alert{\alert{S3 Standard:}} Acceso frecuente.
\item \alert{\alert{S3 Intelligent-Tiering:}} Ajuste automático según el acceso.
\item \alert{\alert{S3 Glacier:}} Archivos de largo plazo.
\end{enumerate}
\end{block}

\begin{block}{Casos de Uso}
\begin{itemize}
\item Almacenamiento multimedia.
\item Copias de seguridad.
\item Análisis de Big Data con Amazon EMR.
\end{itemize}

---
\end{block}
\end{block}

\begin{block}{Google Cloud Storage}
\begin{itemize}
\item Servicio de almacenamiento de objetos de Google Cloud.
\end{itemize}

\begin{block}{Características Clave}
\begin{itemize}
\item Almacenamiento multirregional para redundancia.
\item Consistencia fuerte.
\item Integración con BigQuery.
\end{itemize}
\end{block}

\begin{block}{Clases de Almacenamiento}
\begin{enumerate}
\item \alert{\alert{Standard Storage:}} Datos de acceso frecuente.
\item \alert{\alert{Nearline Storage:}} Acceso menos de una vez al mes.
\item \alert{\alert{Coldline Storage:}} Acceso menos de una vez al año.
\item \alert{\alert{Archive Storage:}} Archivado de datos.
\end{enumerate}
\end{block}

\begin{block}{Casos de Uso}
\begin{itemize}
\item Datos IoT.
\item Logs analíticos.
\item Archivos multimedia distribuidos globalmente.
\end{itemize}

---
\end{block}
\end{block}

\begin{block}{Azure Blob Storage}
\begin{itemize}
\item Sistema de almacenamiento de objetos de Microsoft Azure.
\end{itemize}

\begin{block}{Características Clave}
\begin{itemize}
\item Diseñado para datos no estructurados como imágenes, vídeos y logs.
\item Escalabilidad automática.
\end{itemize}
\end{block}

\begin{block}{Niveles de Almacenamiento}
\begin{enumerate}
\item \alert{\alert{Hot Access Tier:}} Acceso frecuente.
\item \alert{\alert{Cool Access Tier:}} Acceso poco frecuente.
\item \alert{\alert{Archive Tier:}} Archivado a largo plazo.
\end{enumerate}
\end{block}

\begin{block}{Casos de Uso}
\begin{itemize}
\item Copias de seguridad empresariales.
\item Datos analíticos en tiempo real.
\end{itemize}

---
\end{block}
\end{block}
\end{frame}

\begin{frame}[label={sec:org354575d}]{Comparación de Sistemas}
$\backslash$[
\begin{array}{|l|l|l|l|l|l|}
\hline
\textbf{Sistema} & \textbf{Tipo} & \textbf{Escalabilidad} & \textbf{Tolerancia a Fallos} & \textbf{Integración} & \textbf{Costo} \\
\hline
HDFS & Archivos & Alta & Alta (replicación) & Ecosistema Hadoop & Infraestructura propia \\
\hline
Amazon S3 & Objetos & Muy alta & Muy alta & AWS (Lambda, Athena, etc.) & Moderado \\
\hline
Google Cloud Storage & Objetos & Muy alta & Muy alta & GCP (BigQuery, Dataflow) & Moderado \\
\hline
Azure Blob Storage & Objetos & Muy alta & Muy alta & Azure (Synapse, Data Lake) & Moderado \\
\hline
\end{array}
$\backslash$]

---
\end{frame}

\begin{frame}[label={sec:org0b9e43f}]{Conclusión}
Los sistemas de almacenamiento distribuido son esenciales para Big Data.
\begin{itemize}
\item \alert{\alert{On-premises:}} HDFS.
\item \alert{\alert{Nube:}} Amazon S3, Google Cloud Storage, Azure Blob Storage.
\item La elección depende de:
\begin{itemize}
\item Costos.
\item Integración.
\item Requisitos de redundancia y durabilidad.
\end{itemize}
\end{itemize}
\end{frame}
\end{document}
